% !TEX root = ../main.tex
\chapter{関係・前順序・Galois 接続}
\label{chap:relations-preorders-galois}
\BookIndex{かんけい}{関係}
\BookIndex{ぜんじゅんじょ}{前順序}
\BookIndex{Galois connection}{Galois 接続}

\begin{chapterabstract}
本章では、等式だけでは扱いきれない仕様の比較を扱います。
ソフトウェアでは、二つの実装やデータが完全に同じでなくても、一方が他方より詳しい、制約が強い、安全である、抽象的である、と言える場面が多くあります。
このような比較を表す道具が、関係と前順序です。
さらに、詳細な世界と抽象的な世界を往復させるとき、抽象化と具体化の対応を Galois 接続として整理できます。
本章では、公開ログへのマスキングを題材に、情報を落とす処理を「失敗した同型」ではなく「安全な近似」として読む方法を学びます。
\end{chapterabstract}

\begin{learninggoals}
\item 等式ではなく関係で仕様を表す場面を説明できます。
\item 前順序を、反射性と推移性を持つ比較として読めます。
\item 仕様の強弱、情報量、権限、抽象度を順序としてモデル化できます。
\item 単調性を「比較を壊さない変換」として理解できます。
\item Galois 接続を、抽象化と具体化がずれなく対応するための法則として説明できます。
\end{learninggoals}

\section{問題設定}

第18章から第21章までは、主に「同じ結果になること」を扱ってきました。
リファクタリング前後の式が同じであること、関数がすべての入力で同じ出力を返すこと、adapter が業務ロジックと干渉しないことは、等式で表しやすい性質です。

しかし実務には、等式だけでは表しきれない安全性が多くあります。
たとえば、詳細ログから公開ログを作る処理を考えます。
公開ログは、ユーザーID、内部トークン、IPアドレス、細かなエラー理由などを落とすかもしれないため、詳細ログと同じではありません。
したがって、詳細ログと公開ログの間に同型を期待してはいけません。

それでも、公開ログへの変換を正しいと評価したい場合があります。
ここで主張したいのは「元に戻せます」という性質ではありません。
「公開してはいけない情報を含みません」「公開側に残すと決めた項目は保存されています」「詳細なログを抽象化したものとして妥当です」といった性質です。
これらは、等式ではなく、関係や順序の言葉で表すほうが自然です。

図\ref{fig:ch22-equality-vs-order}は、双方向の一致を要求する等式と、向きを持つ比較を区別しています。
情報を落とす変換では round-trip ではなく、保存すべき関係と安全側の順序を選びます。

\FigurePlaceholder{fig-ch22-equality-vs-order.png}{0.90}{等式と順序}{fig:ch22-equality-vs-order}

\section{関係：二つの値を比べる命題}

最初の道具は、関係です。
関係とは、二つの値を受け取って命題を返すものだと考えられます。

たとえば、詳細ログ \texttt{DetailLog} と公開ログ \texttt{PublicLog} の間に、「公開ログの endpoint は詳細ログの endpoint と一致しています」という関係を定義できます。
これは、二つのログが等しいという主張ではありません。
公開ログが詳細ログの一部を正しく反映しているという主張です。

\begin{definitionbox}
型 \(A\) と型 \(B\) の間の関係とは、値 \(a : A\) と \(b : B\) を受け取り、命題を返すものとして読めます。
\[
  R : A \to B \to \mathrm{Prop}
\]
同じ型どうしの関係 \(R : A \to A \to \mathrm{Prop}\) は、等しさ、大小関係、包含関係、互換性、依存関係などを表せます。
\end{definitionbox}

Lean では、関係を \texttt{Prop} を返す関数として書き、その成立を定理として証明します。

\begin{listing}[htbp]
\begin{minted}{lean4}
namespace Chapter22

-- A と B の間の関係は、A の値と B の値から命題を作るもの。
def Rel (A B : Type) := A -> B -> Prop

structure DetailLog where
  userId : Nat
  endpoint : Nat
  token : Nat
deriving Repr

structure PublicLog where
  endpoint : Nat
deriving Repr

-- 詳細ログから公開ログへ落とす。
def maskLog (l : DetailLog) : PublicLog :=
  { endpoint := l.endpoint }

-- 公開ログに残すと決めた部分が保存されている、という関係。
def KeepsPublicPart (d : DetailLog) (p : PublicLog) : Prop :=
  p.endpoint = d.endpoint

theorem maskLog_keeps_public (d : DetailLog) :
    KeepsPublicPart d (maskLog d) := by
  rfl
\end{minted}
\caption{\texttt{KeepsPublicPart}}
\label{lst:ch22-relation}
\end{listing}

この定理は、\texttt{maskLog} が詳細ログを完全に保存することを示しているわけではありません。
\texttt{token} や \texttt{userId} は、公開ログから消えています。
この定理が示すのは、公開ログに残すと決めた \texttt{endpoint} が元の詳細ログと対応していることです。

\begin{warningbox}
情報を落とす処理に対して「元に戻せないから正しくありません」と判断すると、抽象化、匿名化、集計、公開用ビューをすべて不正なものとして扱うことになります。
必要なのは、等式ではなく、その処理に適した関係を定義することです。
\end{warningbox}

\section{前順序：反射性と推移性を持つ比較}

次に、同じ型の値を比較する関係を考えます。
たとえば、ログの詳細度を \texttt{public}、\texttt{internal}、\texttt{secret} の三段階で表すとします。
このとき、\texttt{public} は \texttt{internal} より情報が少なく、\texttt{internal} は \texttt{secret} より情報が少ない、と読めます。

比較が必ずしも数値的な「大小」だけを意味するわけではありません。
前順序は、反射性と推移性を満たす関係です。

\begin{definitionbox}
型 \(A\) 上の関係 \(\sqsubseteq\) が前順序であるとは、次の二つを満たすことです。
\[
\begin{aligned}
&\text{反射性:} && a \sqsubseteq a \\
&\text{推移性:} && a \sqsubseteq b \text{ かつ } b \sqsubseteq c \text{ なら } a \sqsubseteq c
\end{aligned}
\]
反射性は「自分自身と比較できる」ことを意味します。
推移性は「安全な比較を二段つなげても安全である」ことを意味します。
\end{definitionbox}

\begin{table}[htbp]
\centering
\caption{前順序の例}
\label{tab:ch22-preorder-examples}
\begin{tabular}{lll}
\toprule
比較したいもの & \(x \sqsubseteq y\) の読み方 & 実務上の意味 \\
\midrule
ログ詳細度 & \(x\) は \(y\) より情報が少ない & 公開可能性、マスキング \\
仕様 & \(x\) は \(y\) を含意する & \(x\) は \(y\) 以上に強い \\
権限 & \(x\) は \(y\) に含まれる権限である & role / permission 設計 \\
静的解析結果 & \(x\) は \(y\) より精密な近似である & 抽象解釈、警告の安全性 \\
データ形式 & \(x\) は \(y\) へ安全に弱化できる & lossy migration \\
\bottomrule
\end{tabular}
\end{table}

\begin{listing}[htbp]
\begin{minted}{lean4}
-- 前順序の最小構造。
-- le は比較関係、refl は反射性、trans は推移性を表す。
structure PreorderMini (A : Type) where
  le : A -> A -> Prop
  refl : forall a, le a a
  trans : forall {a b c}, le a b -> le b c -> le a c
\end{minted}
\caption{\texttt{PreorderMini}}
\label{lst:ch22-preorder-mini}
\end{listing}

本章では、\(x \sqsubseteq y\) を「\(x\) は \(y\) 以下です」と読みます。
ただし、何が「以下」なのかは場面ごとに決めます。
ログの例では「情報量が少ない」、権限の例では「許可が少ない」が以下です。
仕様の例では、向きを間違えないよう注意が必要です。
一般に、強い仕様は弱い仕様を含意するため、強い仕様を小さい側に置くか大きい側に置くかを、章やプロジェクト内で統一します。

図\ref{fig:ch22-preorder-chain}では、情報量が少ない順に
\texttt{publicLevel}、\texttt{internalLevel}、\texttt{secretLevel} を並べます。

\FigurePlaceholder{fig-ch22-preorder-chain.png}{0.88}{ログ詳細度の順序}{fig:ch22-preorder-chain}

\begin{listing}[htbp]
\begin{minted}{lean4}
inductive Detail where
  | publicLevel
  | internalLevel
  | secretLevel
deriving Repr, DecidableEq

-- public <= internal <= secret と読む。
def detailLe : Detail -> Detail -> Prop
  | .publicLevel, _ => True
  | .internalLevel, .internalLevel => True
  | .internalLevel, .secretLevel => True
  | .secretLevel, .secretLevel => True
  | _, _ => False

theorem detail_refl : forall d : Detail, detailLe d d := by
  intro d
  cases d <;> simp [detailLe]

theorem detail_trans :
    forall {a b c : Detail}, detailLe a b -> detailLe b c -> detailLe a c := by
  intro a b c hab hbc
  cases a <;> cases b <;> cases c <;> simp [detailLe] at *
\end{minted}
\caption{\texttt{detailLe}}
\label{lst:ch22-detail-preorder}
\end{listing}

この順序では、\texttt{public} はどのレベル以下でもあります。
公開情報は、内部情報や秘密情報よりも情報量が少ないためです。
一方、\texttt{secret} は \texttt{public} 以下ではありません。
秘密情報を公開情報として扱うには、抽象化やマスキングが必要です。

\begin{listing}[htbp]
\begin{minted}{lean4}
theorem no_secret_as_public :
    detailLe Detail.secretLevel Detail.publicLevel -> False := by
  intro h
  simp [detailLe] at h
\end{minted}
\caption{\texttt{no\_secret\_as\_public}}
\label{lst:ch22-no-downgrade}
\end{listing}

このコードでは、\texttt{secretLevel}が\texttt{publicLevel}以下であるという仮定を置くと、矛盾することを示しています。
秘密情報を公開情報としてそのまま扱わず、後で出てくる抽象化やマスキングを通す必要がある、と読みます。

\section{仕様の強弱と単調性}

前順序は、データの情報量だけでなく、仕様の強弱にも使えます。
たとえば、ある入力 \texttt{x} に対する述語 \texttt{p x} と \texttt{q x} を考えます。
\texttt{p} が成り立つなら必ず \texttt{q} も成り立つとき、\texttt{p} は \texttt{q} より強い仕様だと言えます。

強い仕様が許す値は少なくなります。
たとえば、「偶数の自然数です」という仕様は「自然数です」という仕様を含意するため、後者より強い仕様です。

\begin{listing}[htbp]
\begin{minted}{lean4}
-- p が q より強いとは、p を満たせば q も満たすということ。
def Stronger {A : Type} (p q : A -> Prop) : Prop :=
  forall x, p x -> q x

theorem stronger_refl {A : Type} (p : A -> Prop) :
    Stronger p p := by
  intro x hx
  exact hx

theorem stronger_trans {A : Type} {p q r : A -> Prop} :
    Stronger p q -> Stronger q r -> Stronger p r := by
  intro hpq hqr x hpx
  exact hqr x (hpq x hpx)
\end{minted}
\caption{\texttt{Stronger}}
\label{lst:ch22-stronger-spec}
\end{listing}

変換が順序を壊さないことも重要です。
この性質を単調性と呼びます。
ログ詳細度の例では、詳細レベルを公開ビューへ抽象化するとき、元の詳細度の比較が抽象後の比較としても成り立つことを求めます。

\begin{definitionbox}
順序 \(\sqsubseteq_A\) を持つ型 \(A\) から、順序 \(\sqsubseteq_B\) を持つ型 \(B\) への関数 \(f : A \to B\) が単調であるとは、次を満たすことです。
\[
  x \sqsubseteq_A y \quad\Rightarrow\quad f(x) \sqsubseteq_B f(y)
\]
単調性は、「比較できていたものを変換後も比較できる」ことを保証します。
\end{definitionbox}

\begin{listing}[htbp]
\begin{minted}{lean4}
def MonotoneMini {A B : Type}
    (leA : A -> A -> Prop) (leB : B -> B -> Prop)
    (f : A -> B) : Prop :=
  forall {x y}, leA x y -> leB (f x) (f y)
\end{minted}
\caption{\texttt{MonotoneMini}}
\label{lst:ch22-monotone}
\end{listing}

このコードでは、二つの順序\texttt{leA}と\texttt{leB}を受け取り、関数\texttt{f}が順序を保つことを\texttt{MonotoneMini}として定義しています。
入力側で比較できる二つの値が、変換後にも出力側の順序で比較できることを要求しています。

\section{抽象化と具体化}

ここから、Galois 接続の準備に入ります。
Galois 接続では、具体的な世界と抽象的な世界という二つの世界を考えます。

詳細ログの例では、具体的な世界は \texttt{Detail} です。
そこでは、\texttt{public}、\texttt{internal}、\texttt{secret} の三段階を区別します。
一方、公開判定では、そこまで細かい区別が不要な場合があります。
公開ビューでは、「公開情報だけ」か「非公開情報を含みうる」かの二段階だけを見ます。

抽象化は情報を落とす関数であると同時に、どの情報を区別しないかを定めます。
本章の \texttt{allow} は、抽象値が許す具体値の最大レベルを返します。

\begin{listing}[htbp]
\begin{minted}{lean4}
inductive PublicView where
  | publicOnly
  | mayContainPrivate
deriving Repr, DecidableEq

-- publicOnly <= mayContainPrivate と読む。
def viewLe : PublicView -> PublicView -> Prop
  | .publicOnly, _ => True
  | .mayContainPrivate, .mayContainPrivate => True
  | .mayContainPrivate, .publicOnly => False

-- 抽象化: 詳細レベルを公開判定用の二段階へ落とす。
def erase : Detail -> PublicView
  | .publicLevel => .publicOnly
  | .internalLevel => .mayContainPrivate
  | .secretLevel => .mayContainPrivate

-- 具体化: 抽象ビューが許す最大の詳細レベルを見る。
def allow : PublicView -> Detail
  | .publicOnly => .publicLevel
  | .mayContainPrivate => .secretLevel
\end{minted}
\caption{\texttt{erase} と \texttt{allow}}
\label{lst:ch22-view}
\end{listing}

\texttt{erase} は、抽象化です。
\texttt{internal} と \texttt{secret} を同じ \texttt{mayContainPrivate} にまとめるため、情報を失います。
したがって、\texttt{erase} は同型ではありません。

一方、\texttt{allow} は具体化です。
\texttt{publicOnly} に対しては、許される具体レベルの最大値を \texttt{public} と定めています。

\texttt{mayContainPrivate} に対しては、許される具体レベルの最大値を \texttt{secret} と定めています。

\begin{listing}[htbp]
\begin{minted}{lean4}
theorem view_refl : forall v : PublicView, viewLe v v := by
  intro v
  cases v <;> simp [viewLe]

theorem erase_monotone :
    MonotoneMini (A := Detail) (B := PublicView) detailLe viewLe erase := by
  intro x y h
  cases x <;> cases y <;> simp [detailLe, viewLe, erase] at *

theorem allow_monotone :
    MonotoneMini (A := PublicView) (B := Detail) viewLe detailLe allow := by
  intro x y h
  cases x <;> cases y <;> simp [viewLe, detailLe, allow] at *
\end{minted}
\caption{\texttt{erase\_monotone} と \texttt{allow\_monotone}}
\label{lst:ch22-erase-monotone}
\end{listing}

単調性により、詳細度の順序を抽象化後にも壊さないことが分かります。
ただし、単調性だけでは、抽象化と具体化が十分に対応しているとは言えません。
そこで、Galois 接続を使います。

図\ref{fig:ch22-abstraction-concretization}では、\texttt{erase} が三段階の詳細度を二段階の公開ビューへまとめ、\texttt{allow} が各ビューの最大詳細度を返します。

\FigurePlaceholder{fig-ch22-abstraction-concretization.png}{0.92}{抽象化と具体化}{fig:ch22-abstraction-concretization}

\section{Galois 接続：抽象化と具体化の対応}

本章では、\(\alpha(c) \leq a \Longleftrightarrow c \leq \gamma(a)\) の向きを採用します。
前順序を圏と見たとき、これは \(\alpha\) を左随伴、\(\gamma\) を右随伴とする向きです~\cite{riehl-category-theory-in-context}。

Galois 接続は、抽象化 \(\alpha\) と具体化 \(\gamma\) が、順序の比較を通じて対応していることを表す法則です。

具体的な世界を \(C\)、抽象的な世界を \(A\) とします。
それぞれに順序 \(\sqsubseteq_C\)、\(\sqsubseteq_A\) があるとします。
抽象化 \(\alpha : C \to A\) と具体化 \(\gamma : A \to C\) が Galois 接続をなすとは、次が成り立つことです。

\[
  \alpha(c) \sqsubseteq_A a
  \quad\Longleftrightarrow\quad
  c \sqsubseteq_C \gamma(a)
\]

左辺は「具体値 \(c\) を抽象化した結果は、抽象値 \(a\) 以下です」と読みます。
右辺は「具体値 \(c\) は、抽象値 \(a\) を具体化した上限以下です」と読みます。
この二つが同値であるとき、抽象側での判定と具体側での判定がずれていないと言えます。

\begin{definitionbox}
本章で使う Galois 接続は、次の形で読めます。
\[
  \text{抽象化したものが抽象仕様を満たす}
  \quad\Longleftrightarrow\quad
  \text{元の具体値が、その抽象仕様の具体化に収まる}
\]
これは、静的解析、権限設計、公開ログ、仕様弱化で現れる「抽象的に安全なら、具体的にも許された範囲にいます」という対応を表します。
\end{definitionbox}

\begin{listing}[htbp]
\begin{minted}{lean4}
structure GaloisConnectionMini (C A : Type)
    (cLe : C -> C -> Prop) (aLe : A -> A -> Prop)
    (alpha : C -> A) (gamma : A -> C) : Prop where
  law : forall c a, Iff (aLe (alpha c) a) (cLe c (gamma a))

-- erase と allow は Galois 接続をなす。
theorem erase_allow_galois :
    GaloisConnectionMini Detail PublicView detailLe viewLe erase allow := by
  constructor
  intro c a
  cases c <;> cases a <;> simp [erase, allow, detailLe, viewLe]
\end{minted}
\caption{\texttt{GaloisConnectionMini}}
\label{lst:ch22-galois-def}
\end{listing}

この証明は、すべての場合を列挙しています。
\texttt{Detail} は三通り、\texttt{PublicView} は二通りなので、合計六通りです。
各ケースで、左辺と右辺が同じ真偽になることを \texttt{simp} で確認しています。

図\ref{fig:ch22-galois-law}の左右は別々の安全条件ではなく、同じ比較を抽象側と具体側で表した同値な命題です。

\FigurePlaceholder{fig-ch22-galois-law.png}{0.94}{Galois 接続の法則}{fig:ch22-galois-law}

\section{Galois 接続から得られる二つの安全性}

Galois 接続があると、抽象化と具体化の往復について、二つの基本的な性質が得られます。

一つ目は、具体値 \(c\) を抽象化してから具体化すると、元の \(c\) を覆う上限になるという性質です。
\[
  c \sqsubseteq_C \gamma(\alpha(c))
\]
これは、「抽象化で情報を落とした後に戻しても、元の値を含む安全な範囲になります」と読めます。

二つ目は、抽象値 \(a\) を具体化してから抽象化すると、元の \(a\) 以下になるという性質です。
\[
  \alpha(\gamma(a)) \sqsubseteq_A a
\]
これは、「抽象仕様から許される最大の具体値を選び、それを再び抽象化しても、元の抽象仕様を超えません」と読めます。

\begin{listing}[htbp]
\begin{minted}{lean4}
theorem gc_unit {C A : Type} {cLe : C -> C -> Prop} {aLe : A -> A -> Prop}
    {alpha : C -> A} {gamma : A -> C}
    (gc : GaloisConnectionMini C A cLe aLe alpha gamma)
    (a_refl : forall a, aLe a a) :
    forall c, cLe c (gamma (alpha c)) := by
  intro c
  exact (gc.law c (alpha c)).mp (a_refl (alpha c))

theorem gc_counit {C A : Type} {cLe : C -> C -> Prop} {aLe : A -> A -> Prop}
    {alpha : C -> A} {gamma : A -> C}
    (gc : GaloisConnectionMini C A cLe aLe alpha gamma)
    (c_refl : forall c, cLe c c) :
    forall a, aLe (alpha (gamma a)) a := by
  intro a
  exact (gc.law (gamma a) a).mpr (c_refl (gamma a))
\end{minted}
\caption{\texttt{gc\_unit} と \texttt{gc\_counit}}
\label{lst:ch22-galois-consequences}
\end{listing}

このコードでは、Galois接続の法則から、具体値を抽象化して戻す方向と、抽象値を具体化して戻す方向の基本的な不等式を取り出しています。
\texttt{gc\_unit}と\texttt{gc\_counit}は、抽象化と具体化が互いにどの範囲を覆うかを確認するための一般形です。

\begin{listing}[htbp]
\begin{minted}{lean4}
theorem detail_covered_by_allowed (d : Detail) :
    detailLe d (allow (erase d)) := by
  exact gc_unit erase_allow_galois view_refl d

theorem allowed_abstraction_is_inside (v : PublicView) :
    viewLe (erase (allow v)) v := by
  exact gc_counit erase_allow_galois detail_refl v
\end{minted}
\caption{\texttt{detail\_covered\_by\_allowed}・\texttt{allowed\_abstraction\_is\_inside}}
\label{lst:ch22-galois-use}
\end{listing}

\texttt{detail\_covered\_by\_allowed} は、詳細レベルを一度抽象化し、その抽象ビューが許す最大の具体レベルへ戻すと、元の詳細レベルを覆うことを述べています。
たとえば、\texttt{internal} を抽象化すると \texttt{mayContainPrivate} になり、それを具体化すると \texttt{secret} になります。
\texttt{internal <= secret} なので、元のレベルを安全に覆っています。

\texttt{allowed\_abstraction\_is\_inside} は、抽象ビューから許される最大の具体レベルを選び、それを再度抽象化しても、元の抽象ビューを超えないことを述べています。
これは、公開ポリシーの上限を越えないことに対応します。

\section{サンプル：詳細ログから公開ログへの安全な抽象化}

ここまでの道具を、公開ログの設計レビューとして読み替えます。

まず、詳細ログには多くの情報が含まれていますが、公開ログにはその一部だけを残します。
この時点で、詳細ログと公開ログは同型ではなく、公開ログから詳細ログを復元できません。

次に、公開ログに残す部分について、関係で仕様を書きます。
たとえば、\texttt{endpoint} は保存されているべきです。
この性質を \texttt{KeepsPublicPart} として表しました。

さらに、ログ全体の安全性を前順序で見ます。
\texttt{public <= internal <= secret} という順序を決めれば、「公開側に出してよい情報か」「内部に留めるべき情報か」を比較できます。

最後に、詳細な分類から公開判定用の分類へ落とす抽象化 \texttt{erase} と、その抽象分類が許す具体的な最大値 \texttt{allow} を定義します。
これらが Galois 接続をなすことを証明すれば、抽象側のポリシー判定と具体側の情報量判定が対応していることを確認できます。

図\ref{fig:ch22-public-log-mask}は、フィールド保存の関係、レベル順序の単調性、抽象化と具体化の対応を別々の証明義務として示します。
Galois 接続だけで、フィールド保存まで自動的に得られるわけではありません。

\FigurePlaceholder{fig-ch22-public-log-mask.png}{0.94}{公開ログの証明義務}{fig:ch22-public-log-mask}

設計レビューでは、次のように分けて確認します。
\begin{itemize}
\item 保存すべき項目は、関係として定義します。
\item 情報量や権限の大小は、前順序として定義します。
\item 変換が順序を壊さないことは、単調性として確認します。
\item 抽象化と具体化の対応は、Galois 接続として確認します。
\end{itemize}
これらを個別に証明して初めて、「公開ログは詳細ログと同じではありませんが、定めた仕様に対して安全な近似です」と言えます。

\section{実務への読み替え}

Galois 接続は、本章の小さな例だけに使うものではありません。
実務では、次のような形で現れます。

\begin{table}[htbp]
\centering
\caption{Galois 接続の応用例}
\label{tab:ch22-galois-applications}
\begin{tabular}{llll}
\toprule
領域 & 具体側 & 抽象側 & 読み方 \\
\midrule
公開ログ & 詳細ログ & 公開ビュー & マスキング後の安全性 \\
権限設計 & 個別 permission & role & role が許す権限範囲 \\
静的解析 & 実行時状態 & 解析上の近似 & 警告が安全側に倒れていること \\
仕様弱化 & 強い仕様 & 公開・互換仕様 & 弱めた仕様で保証する範囲 \\
データ集約 & 生ログ & 集計値 & 集計が元データの上限・下限を保つこと \\
\bottomrule
\end{tabular}
\end{table}

静的解析の例では、具体側は実際のプログラム状態であり、抽象側は「この変数は非負らしいです」「この値は null かもしれません」のような解析結果です。
解析は、実行時状態を完全には再現しません。
実行時に起こりうる状態を安全に覆う近似であることが重要です。

権限設計では、具体側を個別 permission の集合、抽象側を role として読めます。
role を具体化すると、その role が許す permission の集合になります。
個別権限が role の範囲内に収まっているかどうかは、前順序や包含関係で表せます。

lossy migration では、復元ではなく仕様弱化や正規化後の一致として安全性を表せます。

\section{本章のまとめ}

関係は、二つの値が何らかの意味で対応していることを表す命題です。

前順序は、反射性と推移性を持つ比較であり、情報量、権限、仕様の強弱、抽象度を表せます。

単調性は、変換が比較関係を壊さないことを表します。

Galois 接続は、抽象化と具体化が順序を通じて対応していることを表します。

情報を落とす変換は、同型ではなく、安全な近似として仕様化すべき場合があります。

本章では、等式から一歩進み、順序や近似で安全性を扱う方法を学びました。
そこでは、Pullback を「制約を満たすペア」として読み、DB join や複数APIレスポンスの統合へ接続します。
